\section{From mathematical abundance to cumulative research}\label{sec:intro}

Mathematical research produces more than the statements eventually published. It produces representations, intermediate lemmas, tests of plausible routes, failures that reveal necessary hypotheses, and decisions about which question to pursue next. Much of this work remains in notebooks, conversations, drafts, or personal memory. A completed paper selects a coherent argument from that process. This selection is essential for communication, but it can leave later researchers unable to tell why an approach was attempted or which nearby routes have already been tested. When the exploration record is unavailable, both human and AI investigators may spend effort repeating unsuccessful approaches. Such a record is useful only when it preserves the target, assumptions, method, and outcome: an unsuccessful attempt is not by itself a proof that a route is impossible.

AI-assisted mathematical work makes this issue more pressing. Published papers and manuscripts may present a result and its polished proof without preserving the external research history of targets, candidate arguments, tests, and revisions that led to them. An inspectable record of that history could help experts assess how a result was reached and give students material for understanding research decisions, not merely its final proof. The record documents observable research actions and evidence. A large release of manuscripts can contain useful mechanisms even when individual claims require correction or further scrutiny. It can also create a substantial review burden. The quantity of generated text gives little information about mathematical validity, novelty, explanatory value, or usefulness to another researcher. These are different questions, and infrastructure should preserve their differences. The Leiden Declaration on Artificial Intelligence and Mathematics emphasizes human understanding, responsibility, and the needs of reviewers \cite{leiden}. A research system that produces more material without making its assumptions and evidential basis easier to inspect may worsen the problem it aims to solve.

The central proposal of this paper is to make the \emph{research investigation} a reusable mathematical object. Its record should connect a target to the arguments that motivated it, the conditions required for progress, the outcomes of attempted routes, and the changes those outcomes make to subsequent research. By making supported results and carefully scoped unsuccessful attempts available to later investigations, RRE aims to reduce redundant effort and help research accumulate across people and systems. The proposed unit of accumulation is therefore larger than a theorem and more structured than an unedited conversation. It includes a claim, its justification, its scope, its provenance, and the research decisions that made it relevant.

\subsection{Three phases of the Math Instruments program}

Phase I identifies reusable mathematical instruments in research manuscripts \cite{xie1}. An instrument may be a construction, a proof mechanism, or a way to introduce useful structure into a problem. The motivating survey examined an initial collection of 722 AI-written manuscripts and selected 38 problem families for closer study. Its purpose was to identify potentially reusable ideas while retaining uncertainty about validity, antecedents, and novelty.

Phase II reconstructs the reasoning that makes such instruments usable \cite{xie2}. It represents claims and their jointly required premises through source-linked directed hypergraphs, preserving hypotheses, evidence, and transfer obligations. Its six-paper study consists of four OpenAI manuscripts and two retrospective application papers by the author. The study reports 158 arguments and 454 inferences, together with a multiscale visualization. Phase II also proposes a research cycle in which a curated collection supports candidate approaches, hypothesis suggestions, transfer investigations, verification, and the return of outcomes to the collection.

Phase III develops that proposed cycle into \emph{Recursive Reasoning Exploration}. The emphasis moves from extracting an argument to conducting and recording an investigation that can repeatedly extend the argument collection. Building on Phase II, it specifies how these ideas operate together: explicit objectives, generated intermediate targets, recursive investigation, evidence-sensitive graph updates, and declared stopping criteria.

\begin{table}[htbp]
\centering\small
\begin{tabularx}{\linewidth}{@{}p{1.35cm}YYY@{}}
\toprule
Phase & Primary question & Main record & Next use \\
\midrule
I & What mathematical device may be reusable? & Instrument, source, antecedents, caveats & Select mechanisms worth reconstructing \\
II & Why does it work, and under which hypotheses? & Reasoning graph with scope and transfer obligations & Find candidate routes and missing premises \\
III & What should be investigated next, and how does the outcome change what we know? & Evolving reasoning graph and exploration history & Reuse results and failures in subsequent research \\
\bottomrule
\end{tabularx}
\caption{The phases describe a research program. They are not levels of verification: every phase needs its own mathematical review.}
\label{tab:phases}
\end{table}

\subsection{What Recursive Reasoning Exploration means}

RRE combines a collection of reasoning graphs with research objectives, a resource budget, and criteria for accepting results. Objectives may be supplied by researchers or formulated during exploration under a declared research policy. For each investigation, RRE identifies an unresolved obligation, retrieves relevant arguments, proposes a route, and investigates the additional claims that route requires. A result can resolve the target, qualify it, refute a proposed statement, or leave it open. The record is updated accordingly, and the updated record helps determine the next investigation.

The word \emph{recursive} has two meanings. Within an investigation, a proposed argument generates subproblems that are investigated by the same protocol. Across investigations, an outcome changes the available knowledge and can generate a new target. The word \emph{reasoning} identifies the shared object: explicit mathematical arguments that people and AI systems can inspect, revise, and extend.

We distinguish four operating modes according to how objectives are formulated and investigations are conducted (\cref{tab:operating-modes}). The second mode, \emph{Collaborative Reasoning Exploration} (CRE), is also called collaborative RRE. The fourth mode is prospective: an agent would examine the collection, formulate research objectives, propose connections or mathematical constructions, and investigate their consequences without receiving a specific mathematical problem for each investigation. It would still operate under explicit criteria for mathematical value, evidence, and resource use.

\begin{table}[htbp]
\centering\small
\begin{tabularx}{\linewidth}{@{}>{\raggedright\arraybackslash}p{4.1cm}Y@{}}
\toprule
Operating mode & Objective formation and investigation \\
\midrule
Human-directed exploration & Researchers choose and investigate questions, using the graph to organize their work. \\
Collaborative Reasoning Exploration & Humans and AI jointly formulate questions, investigate arguments, and extend the graph, with recorded contributions and handoffs. \\
Goal-directed autonomous exploration & A researcher supplies the overall objective; an AI agent recursively selects and investigates intermediate targets within the declared charter. \\
Open-ended autonomous exploration & An AI agent formulates research objectives from the collection and investigates them under standing evaluation and resource policies. \\
\bottomrule
\end{tabularx}
\caption{Four operating modes of RRE. They describe the allocation of research activities, not levels of mathematical verification. The case study illustrates the third mode.}
\label{tab:operating-modes}
\end{table}

A project can move between modes while preserving one graph and history. Automation does not permit self-certification or relax the rules for admitting a mathematical claim. \emph{Research-strategy improvement} is a separate proposed capability: the system could evaluate and revise its own retrieval, target-selection, or investigation methods across these modes. Adding graph records alone does not establish such an improvement. The case study uses the run identifier RRE005 and illustrates goal-directed autonomous RRE.

Explainability here has an operational meaning. A reader should be able to inspect which premises support a conclusion, why a target was selected, which alternatives were recorded, what was checked, and what remains unresolved. These records are mathematical justifications and externally inspectable research actions. They are not a claim to recover a model's internal computations or a mathematician's unrecorded thought process.

The intended outcome is a shared environment in which human researchers and AI systems can contribute to the same evolving mathematical record. An expert may formulate a question, an AI system may propose a lemma, a computer algebra program may test an identity, and another researcher may discover a counterexample. The record should preserve those different contributions without collapsing them into a single label such as ``AI-generated proof.''

\subsection{Contributions and case-study progression}

This paper makes three connected contributions. First, it specifies RRE as a protocol connecting reasoning graphs, research targets, verification records, and versioned exploration events. Second, it develops a substantive mathematical case study: an exact certificate-degree theorem for all inconsistent two-term tensor decomposition systems, together with nearly sharp bounds for nilpotent matrix pencils. Third, it presents a research agenda for shared human--AI exploration, educational access to research paths, and evaluation of cumulative progress.

The case study begins with two results already in the complexity atlas: a universal lower bound on certificate degree and a matching upper bound when a flattening minor detects excessive rank. The author selected the missing boundary case as a new objective: obtain matching bounds on an infinite family that flattenings miss. The agent first developed exact small certificates, then a uniform argument for nilpotent pencils. The first certificate supplied a concrete base for another target: transfer it to higher tensor orders while charging the contracted modes in the original parameters. Balance was already available from symmetry reduction; the useful additional step was to show that a balanced degree block of size $k$ transfers at cost $n$. The resulting lemma, combined with classical rank-two geometry, gives the exact value $3n$ for every order-$n$ tensor of rank greater than two. These dependencies and the separate historical sequence are the central example of RRE in this paper.

\Cref{sec:case} states the theorem and follows the exploration through six saved graph revisions. \Cref{app:math} supplies the full proofs and the explicit base certificate. \Cref{app:records} documents the event history, exact checks, and replay package; \cref{app:complexity-graphs} explains how the underlying complexity atlas was extracted and how the new component was integrated. The author supplied the objective and milestone; the AI research agent developed the recorded arguments, programs, and subsequent targets. The study documents a mathematical outcome and its research process. Comparative evaluation of the framework is developed in \cref{sec:evaluation}.

\subsection{Intellectual context and related systems}

The idea of making mathematical practice explicit has substantial antecedents. P\'olya's treatment of problem solving emphasizes questions and heuristics that guide a search \cite{polya}. Lakatos examines how conjectures and proofs develop through criticism and revision \cite{lakatos}. RRE's retention of failed routes and scope corrections belongs to this tradition.

Proof planning represents strategies and their applicability conditions to guide theorem proving \cite{bundy}. Modern neural proof search, including HyperTree Proof Search, already uses structured search and learning from earlier proof attempts \cite{htps}. RRE's AND-premise representation and recursive subgoals therefore have clear precedents. Its proposed emphasis is a persistent research record spanning problem formulation, informal and formal arguments, quantitative feasibility, failed transfers, review, and reuse across investigations. A proof-search engine could be one component within that record.

Mathematical knowledge management provides another foundation. MMT and related work address modular theory graphs, large interconnected corpora, and versioned mathematical knowledge \cite{mkm}. The Lean mathematical library demonstrates community development of a reusable formal library \cite{mathlib}. RRE should connect to such systems rather than introduce an incompatible replacement for formal semantics or proof checking. An informal reasoning graph can identify material worth formalizing; it does not inherit kernel verification merely from pointing toward a proof assistant.

Polymath provides a precedent for distributed mathematical collaboration with visible contributions \cite{polymath}. Provenance standards such as W3C PROV describe entities, activities, and agents involved in producing research artifacts \cite{prov}. RRE can reuse these ideas while adding mathematical scope, obligations, and evidence types. A chronology of contributions and a dependency graph answer different questions and should remain linked but distinguishable.

AI-driven scientific exploration also predates this proposal. FunSearch combines program generation, evaluation, and a population of retained programs \cite{funsearch}; The AI Scientist describes an iterative workflow extending from ideas to experiments and reports \cite{aiscientist}. RRE focuses on mathematical argument reuse and the explicit treatment of open obligations, counterexamples, and verification boundaries. The contribution here is their integration around mathematical dependencies, quantitative obligations, and an evolving research record.


The mathematical case study also connects to the geometry and complexity of tensor decomposition. Flattening equations for the second secant variety were established by Landsberg and Manivel and strengthened by Raicu \cite{landsbergmanivel,raicu}; Ballico and Bernardi describe ranks and decompositions on the tangent developable \cite{ballicobernardi}. Aliabadi and Friedland formulate tensor-rank testing through polynomial equations and effective Nullstellensatz linear systems \cite{aliabadifriedland}. Our result determines the exact certificate degree in the original two-term parameters, including tensors of border rank two whose decomposition fibers are empty. Algebraic proof complexity for tensor isomorphism studies a related but different equation system \cite{galesi}. The distinction between those objects guides both the theorem and its graph representation.
